\section{\label{sec:phi4}标量场理论的梯度流}

我们在四维欧氏标量场论中工作，该理论具有四次自相互作用与裸质量 $m$，由作用量
\begin{equation}
S_\phi[\phi] = \frac{1}{2}\int \mathrm{d}^4x\, \left[(\partial_\nu
\phi)^2 + m^2\phi^2 + \frac{\lambda}{12}\phi^4\right].
\end{equation}
定义。

为研究 sOPE，我们引入标量梯度流，它由流时间演化方程
\begin{equation}
\frac{\partial \rho(\tau,x)}{\partial \tau}= \partial^2 \rho,
\end{equation}
定义，其中 $\rho$ 是非零流时间 $\tau$ 处的标量场，$\partial^2$
是欧氏四维拉普拉斯算子。注意流时间的量纲为 $[\tau] = [x]^2$。

施加狄利克雷边界条件 $\rho(0,x) = \phi(x)$ 后，
我们可以把流时间方程的精确解写作
\begin{equation}\label{eq:rhoexp}
\rho(\tau,x) = e^{\tau\partial^2}\phi(x),
\end{equation}
或者，在动量表示中写作
\begin{equation}
\tilde{\rho}(\tau,p) = e^{-\tau p^2}\tilde{\phi}(p).
\end{equation}

完全解为
\begin{align}
\rho(\tau,x) = {} & \int\mathrm{d}^4y\,\int \frac{\mathrm{d}^4p}{(2\pi)^4}
\,e^{ip\cdot(x-y)} e^{-\tau p^2} \phi(y) \nonumber\\
= {} &  \frac{1}{16\pi^2\tau^2}
\int\mathrm{d}^4y\,e^{-(x-y)^2/4\tau}\phi(y),
\end{align}
它明确展示了梯度流的``涂抹''效应：流时间指数式地压低紫外模式。
我们用方均根涂抹长度 $s_{\mathrm{rms}} = \sqrt{8\tau}$ 来参数化涂抹半径。

流时间为 $\tau_1$ 与 $\tau_2$ 的两个场的（欧氏）涂抹标量传播子由下式给出：
\begin{equation}
\widetilde{G}_\rho(\tau_1,\tau_2,k) = \frac{e^{-(\tau_1+\tau_2) k^2}}{k^2+m^2}.
\end{equation}
流时间演化是经典的，因此任何相互作用都必须发生在零流时间处。四点顶点相应的费曼规则就是标准的
（欧氏）四点顶点，$V^{(4)} = -\lambda/24$。
